\section*{Chapter \ref{ch:b1:money-and-inflation} --- Money and Inflation}
% Numbers recomputed in code/bachelor-1/23-money-and-inflation/exercises/solutions.py (tested).

\begin{solution}{exo:b1:money-and-inflation:qt}
$\pi = 8 - 3 = 5\%$.
\end{solution}

\begin{solution}{exo:b1:money-and-inflation:read}
$PY = MV$, so with the same $M$ nominal GDP would have been $1.782/1.182 = 1.508$ times
its actual value.
\end{solution}

\begin{solution}{exo:b1:money-and-inflation:seign}
$0.2 \times 0.1 = 2\%$ of GDP.
\end{solution}

\begin{solution}{exo:b1:money-and-inflation:fisher-ex}
It rises by 3 points.
\end{solution}

\begin{solution}{exo:b1:money-and-inflation:elast}
$\ln V = -a + 4i$ (with income elasticity one), so $d\ln V/d\ln i = 4i = 0.2$.
\end{solution}

\begin{solution}{exo:b1:money-and-inflation:max}
$\pi^* = 1/2$, and \omterm{def:b1:money-and-inflation:seigniorage}{seigniorage} $\pi e^{-2\pi}Y$ is then $Y/(2e) = 0.1839\,Y$.
\end{solution}

\begin{solution}{exo:b1:money-and-inflation:bolivia}
$12.1/27.1 = 0.446$: real balances relative to GDP fell by more than half.
A hundredfold rise in inflation raised the tax rate far more than it shrank
the base, but the base shrank fast.
\end{solution}

\begin{solution}{exo:b1:money-and-inflation:velfall}
$\pi = g_M + g_V - g_Y = 6 - 1 - 2 = 3\%$: falling \omterm{def:b1:money-and-inflation:velocity}{velocity} means people hold more money
per unit of income, which absorbs part of the money growth.
\end{solution}

\begin{solution}{exo:b1:money-and-inflation:fisher-slope}
Since 1960, inflation passed almost one for one into short-term rates, as
the Fisher effect predicts for expected inflation. Over all decades, much
inflation --- under the gold standard, in wars --- was not expected to last,
so it did not enter nominal rates: the Fisher effect concerns expected
inflation.
\end{solution}

\begin{solution}{exo:b1:money-and-inflation:derive}
$S'(\pi) = ke^{-\alpha\pi}(1 - \alpha\pi)$ is zero at $\pi = 1/\alpha$, positive before and negative after;
$S(1/\alpha) = k/(\alpha e)$. With growth $g$, money must grow at $\pi + g$ to keep $m$ constant,
so $S = (\pi + g)ke^{-\alpha\pi}$, with $S' = ke^{-\alpha\pi}[1 - \alpha(\pi + g)]$: the maximum is at $\pi = 1/\alpha - g$.
\end{solution}

\begin{solution}{exo:b1:money-and-inflation:rule}
$\pi_t = g_M + \bar g_V + \varepsilon_t - g_Y$: mean $g_M + \bar g_V - g_Y$, variance $\sigma^2$. A bank that sets
$g_{M,t} = \bar g_M - \varepsilon_t$ removes the variance entirely: targeting inflation does what
targeting money cannot when \omterm{def:b1:money-and-inflation:velocity}{velocity} moves.
\end{solution}

\begin{solution}{exo:b1:money-and-inflation:refute}
Years or countries with low inflation in which inflation moves with
energy prices, wages or exchange rates while money growth does not, and
money growth moves without inflation following. The long-run link survives
because over decades and across countries, differences in money growth
are large --- from a few per cent to hundreds --- and dwarf the movements in
\omterm{def:b1:money-and-inflation:velocity}{velocity}.
\end{solution}

\begin{solution}{pb:b1:money-and-inflation:kestria}
\textbf{1.} $m = 0.25e^{-0.125} = 0.2206$; $V = 1/m = 4.533$.

\textbf{2.} $m = 0.1947$; $V = 5.136$.

\textbf{3.} $2.5$; the elasticity is $-2.5 \times 0.05 = -0.125$ ($0.125$ in absolute value).

\textbf{4.} $0.05 \times 0.2206 = 1.103\%$ of GDP.

\textbf{5.} Inflation is the cost of holding money: people hold less of it
per unit of income, and spend it faster.

\textbf{6.} $0.25\pi e^{-2.5\pi} = 0.02$.

\textbf{7.} $\pi = 0.1037$ and $\pi = 1.0171$.

\textbf{8.} $m = 0.1929$ and $0.0197$.

\textbf{9.} The low one: starting from low inflation, revenue rises with $\pi$
until the deficit is covered.

\textbf{10.} If people expect the high rate, they hold little money, and
covering the deficit then requires the high rate: expectations can make the
bad \omterm{def:b1:models-and-data:equilibrium}{equilibrium} happen.

\textbf{11.} $\pi^* = 1/2.5 = 0.4$.

\textbf{12.} $e^{0.4} - 1 = 49.18\%$ a year.

\textbf{13.} $0.25/(2.5e) = 3.679\%$ of GDP, with $m = 0.092$.

\textbf{14.} $3.347\%$ and $0.882\%$.

\textbf{15.} No steady inflation rate yields it: the government prints ever
faster, real balances collapse and inflation explodes --- a \omterm{def:b1:money-and-inflation:hyperinflation}{hyperinflation}.

\textbf{16.} \omterm{def:b1:money-and-inflation:seigniorage}{Seigniorage} is a share of GDP because real balances are; GDP
scales both.

\textbf{17.} $m = 0.25e^{-3.75} = 0.00588$, about $0.6\%$ of GDP; prices rise by
$e^{1.5} - 1 = 348.2\%$ a year.

\textbf{18.} $m = 0.25e^{-0.05} = 0.2378$.

\textbf{19.} About $40.4$ times.

\textbf{20.} $23.78\% - 0.59\%$, that is $23.19\%$ of GDP, issued once.

\textbf{21.} Without it, the government would have to keep printing to cover
the deficit, and people would not believe inflation would stay low; real
balances would not rise.

\textbf{22.} Bolivia's broad money went from $10.7\%$ of GDP in 1986 to $46\%$ in
1994, with inflation near $10$ to $20\%$: the same rebuilding of real
balances, spread over years.

\textbf{23.} A deficit that returns, or a central bank seen to finance it;
real balances would stay low, and the new money issued would raise prices.

\textbf{24.} The largest deficit Kestria can finance by printing is $3.7\%$ of
GDP ($3.679\%$, at an inflation rate of $0.4$).
\end{solution}
